% English translation of the French appendix in Backmatter.tex.
\chapter[English Translation of the Appendix]{Cohomology with Proper Support and Construction of the Functor $f^!$}
\markboth{Appendix: English Translation}{Appendix: English Translation}
\begingroup\setlength{\emergencystretch}{2em}
\noindent\textit{English translation of the preceding French appendix. The original statements and provisional arguments are retained; notes marked editorial translate the editorial notes in the French text.}\par\medskip
by P.\ Deligne\footnote{This is an expanded version of a letter from P.\ Deligne to R.\ Hartshorne (letter of 3 March 1966). The footnotes were added by the copyist.}

Verdier has shown that, in the topological case, the formalism of Poincar\'e duality reduces to problems that are local on the source (see \hyperref[appen:bib1]{[1]}). To transpose his construction to the setting of schemes, one needs a theory of cohomology ``with proper support'' for coherent sheaves.

Unless explicitly stated otherwise, all preschemes considered are noetherian and all presheaves are quasi-coherent.

\section*{No.\ 1. The calculus of pro-objects.}\label{appen:sec1}

\begin{bookproposition}{1}\label{appen:prop1}
Let $\mathcal{C}$ be a $\mathbb{U}$-category\footnote{$\mathbb{U}$ denotes a universe fixed throughout.} in which finite colimits exist. Let $h$ be a functor $\mathcal{C}^{\circ}\to(\text{sets})$. The following conditions are equivalent:
\begin{enumerate}
\item[(i)] $h$ is a colimit of representable functors indexed by a small\footnote{``Small'' means ``$\in\mathbb{U}$''.} directed partially ordered set.
\item[(ii)] $h$ is a colimit of representable functors indexed by a small filtered category.
\item[(iii)] $h$ takes finite $\varinjlim$ to finite $\varprojlim$, and there is a small set of objects such that every element of an $h(X)$ factors through one of them.
\end{enumerate}
\end{bookproposition}

The implications (i) $\Rightarrow$ (ii) $\Rightarrow$ (iii) are trivial, and (iii) $\Rightarrow$ (ii) is standard ($h$ is the colimit of the $h_F$ over the category of pairs $(F,\alpha)$ with $\alpha\in h(F)$ and $F$ in the subcategory of $\mathcal C$ closed under finite $\varinjlim$ generated by the given small set). It remains to prove (ii) $\Rightarrow$ (i). If $\mathcal L$ and $\mathcal M$ are two filtered categories, a functor $F:\mathcal L\to\mathcal M$ is called cofinal if $\varinjlim_{L\in\mathcal L}h_{F(L)}$ is the terminal functor. Then for every $G:\mathcal M\to\mathcal N$, one has $\varinjlim_{\mathcal M}G=\varinjlim_{\mathcal L}(G\circ F)$.

\par\medskip\noindent\textbf{Lemma.} For every small filtered category $\mathcal L$, there is a cofinal functor from a small directed partially ordered set to $\mathcal L$.

The first projection $\mathcal L\times\mathbb N\longrightarrow\mathcal L$ is a cofinal functor (with the natural order on $\mathbb N$); this allows us to reduce to the case in which $\operatorname{Ob}\mathcal L$, preordered by $\Hom(X,Y)\ne\emptyset$, has no greatest element. Order by inclusion the set $E$ of finite subcategories of $\mathcal L$ having a unique terminal object (finite means that the set of arrows is finite). Define a functor from $E$ to $\mathcal L$ by assigning to each such category its terminal object. Under the hypotheses made, $E$ is directed and this functor is cofinal.

The functors satisfying the equivalent conditions of \hyperref[appen:prop1]{Proposition 1} are the \emph{Ind-objects} of $\mathcal C$. They form a full subcategory $\operatorname{Ind}\mathcal C$ of $\operatorname{Fun}(\mathcal C^{\mathrm{op}},\mathrm{Ens})$, closed under colimits. The assignment $\mathcal C\longmapsto\operatorname{Ind}\mathcal C$ is a functor $(\text{cat})\to(\text{cat})$. If filtered colimits exist in $\mathcal C$, then for every $h\in\operatorname{Ind}\mathcal C$ the functor $X\longmapsto\Hom(h,h_X)$ is corepresentable, giving a functor, denoted $\varinjlim$, from $\operatorname{Ind}\mathcal C$ to $\mathcal C$. In particular, there is always a functor $\operatorname{Ind}\operatorname{Ind}\mathcal C\longrightarrow\operatorname{Ind}\mathcal C$, and every functor $\mathcal C\longrightarrow\operatorname{Ind}\mathcal D$ extends to a functor $\operatorname{Ind}\mathcal C\longrightarrow\operatorname{Ind}\mathcal D$. In $\operatorname{Ind}\mathcal C$, filtered colimits are exact and will be denoted by ``$\varinjlim$''; in particular, identifying $\mathcal C$ with a full subcategory of $\operatorname{Ind}\mathcal C$, if $(X_i)_{i\in I}$ is a filtered direct system in $\mathcal C$, then $\varinjlim_{i\in I}X_i=\varinjlim\bigl(\text{``}\varinjlim_{i\in I}\text{''}X_i\bigr)$.

For a functor on $\operatorname{Ind}\mathcal C$ to be representable, it is necessary and sufficient that it take finite $\varinjlim$ to finite $\varprojlim$, filtered $\varinjlim$ to filtered $\varprojlim$, and that its restriction to $\mathcal C$ satisfy the smallness condition of \hyperref[appen:prop1]{Proposition 1 (iii)} (indeed, \hyperref[appen:prop1]{Proposition 1} then shows that its restriction to $\mathcal C$ is an Ind-object, with which it agrees everywhere by the condition on limits).

The preceding discussion, except for \hyperref[appen:prop1]{Proposition 1 (iii)} and the preceding assertion, remains valid if one uses sequential limits throughout (or countable limits, which amounts to the same thing).

By duality one defines the category $\operatorname{pro}\mathcal C$ of pro-objects (a full subcategory of $\operatorname{Fun}(\mathcal C,\mathrm{Ens})^{\mathrm{op}}$).

\begin{bookproposition}{2}\label{appen:prop2}
Let $X$ be a quasi-compact, quasi-separated prescheme (not necessarily noetherian). The category of quasi-coherent sheaves on $X$ is equivalent to the category of Ind-objects of the category of finitely presented quasi-coherent sheaves on $X$.
\end{bookproposition}

The map is $\text{``}\varinjlim\text{''}\mathcal F_i\longmapsto\varinjlim_i\mathcal F_i$. If $\mathcal F$ is f.p.\ (finitely presented), then for every filtered direct system $(\mathcal G_i)_{i\in I}$ one has
\[
\Hom_X\bigl(\mathcal F,\varinjlim_i\mathcal G_i\bigr)
\cong\varinjlim_i\Hom_X(\mathcal F,\mathcal G_i)
\] (by a finite gluing argument, which commutes with $\varinjlim$, one reduces to the affine case). The preceding functor is therefore fully faithful. Its image is closed under filtered colimits and finite direct sums. It suffices to prove that for every $x\in X$, every quasi-coherent sheaf $\mathcal F$ on $X$, and every $s\in\mathcal F_x$, there is a finitely presented sheaf $\mathcal G$ and a map $\mathcal G\to\mathcal F$ whose image contains $s$: $\mathcal F$ will be a colimit of images of finitely presented sheaves, each of which is a quotient of a finitely presented sheaf by a colimit of subsheaves of finite type.

Suppose that on a quasi-compact neighborhood $U$ of $x$ there is a map $f:\mathcal C\longrightarrow\mathcal F$, with $\mathcal C$ finitely presented and $s$ in the image. It suffices to extend $\mathcal C$ and $f$ to all of $X$; proceeding step by step, it suffices to know how to extend to $U\cup V$ when $V$ is affine, which amounts to extending from $U\cap V$ to $V$:

\par\medskip\phantomsection\label{appen:extension-lemma}\noindent\textbf{Lemma.}
Let $\mathcal C$ be a finitely presented sheaf on a quasi-compact open subset $U$ of an affine scheme $X$, let $\mathcal F$ be a sheaf on $X$, and let $f$ be a map from $\mathcal C$ to $\mathcal F|_U$. It is possible to extend $\mathcal C$ and $f$ to $\overline{\mathcal C}$ and $\overline f$ defined on $X$ (with $\overline{\mathcal C}$ finitely presented).
\par\medskip

Let $\iota$ be the inclusion of $U$ in $X$, and let $\mathcal C_1$ be the fibre product of $\iota_*\mathcal C$ and $\mathcal F$ over $\iota_*\iota^*\mathcal F$. Then $\mathcal C_1$ extends $\mathcal C$ and maps to $\mathcal F$. If one expresses $\mathcal C_1$ as the colimit of its subsheaves of finite type, then, since $U$ is quasi-compact and $\mathcal C$ is finitely presented, one of them already extends $\mathcal C$. Expressing this last sheaf as a quotient of $\mathcal O^n$ by a colimit of subsheaves of finite type of $\mathcal O^n$, one sees in the same way that $\mathcal C_1$ can be replaced by a finitely presented sheaf.

\begin{bookcorollary}{1}\label{appen:cor1}
For a contravariant additive functor $F:(\text{q coh})\to\Ab$ to be representable, it is necessary and sufficient that it be left exact and take filtered $\varinjlim$ to filtered $\varprojlim$.
\end{bookcorollary}

\begin{bookcorollary}{2}\label{appen:cor2}
Every finitely presented sheaf defined on a quasi-compact open subset $U$ of $X$ extends to a finitely presented sheaf on $X$.
\end{bookcorollary}

Let $\mathcal A$ be an abelian category. We shall work primarily in $\operatorname{pro}\DD^b\mathcal A$, the full subcategory of $\operatorname{pro}\DD(\mathcal A)$ formed by the ``$\varprojlim$'' $K_i$ whose cohomology is uniformly bounded with respect to $i$. The functor $H^p$ extends to a functor $\operatorname{pro}\DD(\mathcal A)\longrightarrow\operatorname{pro}\mathcal A$.

\begin{bookproposition}{3}\label{appen:prop3}
The functors $H^p$ form a conservative system\footnote{That is, if $u$ is a morphism such that the $H^p(u)$ are invertible, then $u$ is invertible.} of functors on $\operatorname{pro}\DD^b\mathcal A$.
\end{bookproposition}

Let $f:K\to L$ be a morphism of $\operatorname{pro}\DD^b\mathcal A$ such that the $H^p(K)\to H^p(L)$ are isomorphisms (in $\operatorname{pro}\mathcal A$). We must verify that for every $M$ in $\DD^b(\mathcal A)$, $\Hom(K,M)\longleftarrow\Hom(L,M)$ is an isomorphism. Write $K=\text{``$\varprojlim$''}K_{\kappa}$. There is a convergent spectral sequence\footnote{Editorial index correction: the original text has $H^qK_{\kappa}$. For the abutment in degree $p+q$, one needs $H^{-q}K_{\kappa}$.}
\[
\Ext^p(H^{q}K_{\kappa},M)\Longrightarrow\Ext^{p+q}(K_{\kappa},M),
\]
whence, on passing to the limit,
\[
\varinjlim_{\kappa}\Ext^p(H^{q}K_{\kappa},M)
\Longrightarrow\varinjlim_{\kappa}\Ext^{p+q}(K_{\kappa},M),
\]
which is still convergent because $K$ is uniformly bounded. Using the fact that every
functor (here $\Ext$) extends to pro-objects, together with the functor
$\varinjlim:\operatorname{Ind}\Ab\to\Ab$, we can rewrite this as
\[
\varinjlim\Ext^p(H^{-q}K,M)
\Longrightarrow\varinjlim\Ext^{p+q}(K,M)
=\Hom(K[-p-q],M).
\]
This spectral sequence also remains valid for $M\in\DD(\mathcal A)$. The proposition follows immediately from the existence of these spectral sequences.

\section*{No.\ 2. Fundamental lemmas.}\label{appen:sec2}

Recall that the preschemes under consideration are noetherian.

\begin{bookproposition}[Duality theorem for an open immersion]{4}\label{appen:prop4}
Let $U$ be an open subset of $X$, $j$ the inclusion map, $\mathcal J$ an ideal defining the complementary closed subset, $\mathcal F$ a coherent sheaf on $U$, and $\mathcal G$ a quasi-coherent sheaf on $X$. Let $\overline{\mathcal F}$ be a coherent extension of $\mathcal F$. One has\footnote{This formula is of course well known; see, for example, P.\ Gabriel's thesis.}
\[
\Hom_{U}(\mathcal{F}, j^{*}\mathcal{G}) = \Hom_{X}(\text{``}\varprojlim_n\text{''} \mathcal{J}^{n}\overline{\mathcal{F}}, \mathcal{G}).
\]
\end{bookproposition}

The map $\varinjlim\Hom_X(\mathcal J^n\overline{\mathcal F},\mathcal G)\longrightarrow\Hom_U(\mathcal F,\mathcal G)$ is injective: if the image under $f$ of $\mathcal J^n\overline{\mathcal F}$ in $\mathcal G$ is supported in the complement of $U$, it is annihilated by a power of $\mathcal J$, say $\mathcal J^k$, and the image of $\mathcal J^{n+k}\overline{\mathcal F}$ is zero.

Now suppose that $X$ is affine, and let $f\in\Hom_U(\mathcal F,\mathcal G)$. For $k$ sufficiently large, every section $s$ of $\mathcal J^k\overline{\mathcal F}$ over $X$ has an image (over $U$) that extends to a section of $\mathcal G$ over $X$. Replacing $\overline{\mathcal F}$ by $\mathcal J^k\overline{\mathcal F}$, we may suppose that we have a diagram
\[
\begin{tikzcd}
0 \ar[r] & R \ar[r] \ar[d,dashed] & \mathcal{O}^{n} \ar[r] \ar[d,equals] & \overline{\mathcal F} \ar[r] \ar[d,dashed] & 0 \\
0 \ar[r] & R_1 \ar[r] & \mathcal{O}^{n} \ar[r] & \mathcal G
\end{tikzcd}
\]
in which the dotted arrows are defined over $U$. For $\ell$ sufficiently large, by Krull and the fact that $R/R\cap R_1$ is supported outside $U$, one has $R\cap\mathcal J^{\ell}\cdot\mathcal O^n\subset R_1$, which allows $f$ to be extended to $\overline f:\mathcal J^{\ell}\overline{\mathcal F}\to\mathcal G$.

In the general case, let the $U_i$ be a finite affine cover of $X$, and let the $U_{ijk}$ be finite affine covers of $U_{ij}=U_i\cap U_j$. Filtered $\varinjlim$ commute with finite products. Set, with restrictions to the relevant open subsets understood,
\[
\begin{aligned}
L(V)&=\Hom_{U\cap V}(\mathcal F,j^*\mathcal G),\\
P(V)&=\varinjlim_n\Hom_V(\mathcal J^n\overline{\mathcal F},\mathcal G).
\end{aligned}
\]
We obtain the following diagram with exact rows:
\[
\begin{tikzcd}[column sep=large,row sep=large]
0 \ar[r] & L(X) \ar[r]
 & \prod_i L(U_i) \ar[r,shift left=.5ex] \ar[r,shift right=.5ex]
 & \prod_{i,j,k}L(U_{ijk}) \\
0 \ar[r] & P(X) \ar[r] \ar[u]
 & \prod_i P(U_i) \ar[r,shift left=.5ex] \ar[r,shift right=.5ex] \ar[u,"\sim"']
 & \prod_{i,j,k}P(U_{ijk}) \ar[u,"\sim"']
\end{tikzcd}
\]
which completes the proof.

Proposition~\ref{appen:prop4} gives an explicit formula for the functor
\[
(\text{coh on }U)\to\operatorname{pro}(\text{coh on }X)
\]
that is ``left adjoint'' to $j^*$. We shall denote this functor by $j_!$
(``extension by zero'')\footnote{This terminology and notation clearly conflict with those generally accepted (Godement's book, SGA 1962, \dots).};
it follows from Krull that it is exact; this also follows from the fact that,
since $X$ is noetherian, an injective object in the category of
quasi-coherent sheaves on $X$, restricted to $U$, remains quasi-coherent and injective on $U$.

\begin{bookproposition}[Independence of the compactification]{5}\label{appen:prop5}
Let $f:X\to Y$ be a proper morphism inducing an isomorphism
between the open subset $U$ of $Y$ and $f^{-1}(U)$:
\[
\begin{tikzcd}
U \ar[r,hook] \ar[d,equals] & X \ar[d,"f"] \\
U \ar[r,hook] & Y.
\end{tikzcd}
\]
Let $\mathcal J$ be a sheaf of ideals defining $Y\setminus U$,
and let $\mathcal J'=\mathcal J\mathcal O_X$ be the induced ideal on $X$.
If $\mathcal F$ is coherent on $X$, one has
\begin{enumerate}
\item[(i)] if $k>0$, the inverse system
$\RR^k f_*(\mathcal J'{}^n\mathcal F)$ is essentially zero
(that is, it defines the zero pro-object);
\item[(ii)] if $k=0$, then for $n$ sufficiently large,
\[
f_*(\mathcal J'{}^{n+1}\mathcal F)
=\mathcal J\cdot f_*(\mathcal J'{}^n\mathcal F).
\]
\end{enumerate}
\end{bookproposition}

One has
\[
\bigoplus_{n\ge0}\RR^k f_*(\mathcal J'{}^n\mathcal F)
\quad\text{is a }\bigoplus_{n\ge0}\mathcal J^n\text{-module of finite type}
\quad\text{\cite[\nopp III]{EGA}}.
\]
Thus, for $n$ sufficiently large,
\[
\mathcal J\otimes\RR^k f_*(\mathcal J'{}^n\mathcal F)
\longrightarrow\RR^k f_*(\mathcal J'{}^{n+1}\mathcal F)
\]
is surjective; (ii) follows. If $k>0$, $\RR^k f_*$ vanishes on $U$;
there exists $N$ such that $\mathcal J^N\RR^k f_*(\mathcal J'{}^n\mathcal F)$
is zero (for every $n$).
To shorten the diagram, put $M_n=\RR^k f_*(\mathcal J'{}^n\mathcal F)$.
The diagram
\[
\begin{tikzcd}[row sep=large,column sep=large]
\mathcal J^p\otimes M_n \ar[r] \ar[d]
 & M_{n+p} \ar[r] \ar[d] & 0\quad(n\ge n_0) \\
\mathcal J^p\cdot M_n \ar[r] & M_n
\end{tikzcd}
\]
shows that for $N$ sufficiently large and $k>0$,
\[
\RR^k f_*(\mathcal J'{}^{n+N}\mathcal F)
\longrightarrow\RR^k f_*(\mathcal J'{}^n\mathcal F)
\]
is zero, which is (i).

\section*{No.\ 3. Definition of $\RR f_!$.}\label{appen:sec3}

A morphism $f$ (respectively, a pair of composable morphisms, a triple, \dots) will be called \emph{compactifiable} if one can find, respectively, commutative diagrams
\[
\begin{tikzcd}[row sep=large,column sep=large]
X \ar[r,hook] \ar[d,"f"'] & \overline X \ar[dl] \\
S
\end{tikzcd}
\qquad
\begin{tikzcd}[row sep=large,column sep=large]
X \ar[r,hook] \ar[d,"g"'] & \overline X \ar[r,hook] \ar[dl]
 & \overline{\overline X} \ar[dl] \\
Y \ar[r,hook] \ar[d,"f"'] & \overline Y \ar[dl] \\
S
\end{tikzcd}
\]
\[
\begin{tikzcd}[row sep=large,column sep=large]
X \ar[r,hook] \ar[d,"h"'] & \overline X \ar[r,hook] \ar[dl]
 & \overline{\overline X} \ar[r,hook] \ar[dl]
 & \overline{\overline{\overline X}} \ar[dl] \\
Y \ar[r,hook] \ar[d,"g"'] & \overline Y \ar[r,hook] \ar[dl]
 & \overline{\overline Y} \ar[dl] \\
Z \ar[r,hook] \ar[d,"f"'] & \overline Z \ar[dl] \\
S
\end{tikzcd}
\]
in which the horizontal arrows are open immersions and the oblique arrows are proper. Every compactifiable morphism is separated and of finite type, and Nagata asserts in \hyperref[appen:bib2]{[2]} that his proof establishes the converse for integral noetherian schemes, but his hypotheses on the base are unclear.\footnote{Mumford is said to have verified, using Nagata's proof, that every separated morphism of finite type between noetherian schemes is compactifiable.}

For every compactifiable $f$, we propose to define
\[
\RR f_!:\operatorname{pro}\DD_{\operatorname{coh}}^b(X)\longrightarrow\operatorname{pro}\DD_{\operatorname{coh}}^b(S).
\] Recall that $\DD_{\operatorname{coh}}^b(S)$, the bounded derived category of the category of coherent sheaves on $S$, is a full subcategory of the derived category of all sheaves (not necessarily quasi-coherent) on $S$, with essential image consisting of the complexes with bounded coherent cohomology.

If $f$ is proper, we have $\RR f_*:\DD_{\operatorname{coh}}^b(X)\longrightarrow\DD_{\operatorname{coh}}^b(S)$, of finite cohomological dimension, which induces $\RR f_!:\operatorname{pro}\DD_{\operatorname{coh}}^b(X)\longrightarrow\operatorname{pro}\DD_{\operatorname{coh}}^b(S)$. If $f$ is an open immersion, then $f_!:(\text{coh on }X)\longrightarrow\operatorname{pro}(\text{coh on }S)$ induces $\DD^b(\text{coh on }X)\longrightarrow\DD^b\operatorname{pro}(\text{coh on }S)$, which maps to $\operatorname{pro}\DD_{\operatorname{coh}}^b(S)$, and this map extends to $\RR f_!:\operatorname{pro}\DD_{\operatorname{coh}}^b(X)\longrightarrow\operatorname{pro}\DD_{\operatorname{coh}}^b(S)$. If the coherent complex $K$ on $X$ is extended to $\overline K$ on $S$, its image is $\text{``}\varprojlim_n\text{''}\mathcal J^n\overline K$, where $\mathcal J$ defines $S\setminus X$.

For a composite $f=gh$ ($g$ proper, $h$ an open immersion), take $\RR f_!=\RR g_!\RR h_!$. We must verify:

\subsection*{Independence of the compactification.}

Consider a commutative diagram
\[
\begin{tikzcd}[row sep=large,column sep=large]
X \ar[r,hook,"j'"] \ar[dr,hook,"j''"'] \ar[dd,"f"']
 & X' \ar[d,"g"] \\
 & X'' \ar[dl,"\overline f"] \\
S
\end{tikzcd}
\]
in which $j'$ and $j''$ are open immersions, and $\overline f$ and $g$
are proper maps. $j'X=g^{-1}j''X\cap j'X$. We must prove that
\[
\RR\overline f_*\RR j''_!
=\RR(\overline f g)_*\RR j'_!.
\]
Since $\RR(\overline f g)_*=\RR\overline f_*\RR g_*$, it suffices to prove
\[
\RR j''_!\cong\RR g_*\RR j'_!.
\]
We reduce to the case in which $X$ is dense in $X'$, so as to meet the
hypotheses of Proposition~\ref{appen:prop5}.
Let $K$ be a bounded coherent complex on $X$, extended to $\overline K$
on $X'$. The complex $g_*\overline K$ extends $K$ to $X''$, and there is a map
\begin{equation}\tag{1}\label{appen:eq:independence}
g_*(\mathcal J^n\overline K)\longrightarrow\RR g_*(\mathcal J^n\overline K).
\end{equation}
Here $\mathcal J$ defines $X'\setminus X$.\footnote{Editorial clarification: one may choose $\mathcal J$ to be induced by an ideal defining $X''\setminus X$; this is the choice used to apply Proposition~\ref{appen:prop5}.}
Part (ii) of that proposition gives, in the category of pro-objects,
\[
j''_!K\cong\text{``}\varprojlim_n\text{''}\,g_*(\mathcal J^n\overline K).
\]
On passing to the formal inverse limit, we obtain the spectral sequence
\[
\text{``}\varprojlim_n\text{''}\,\RR^p g_*(\mathcal J^n\overline K^q)
\Longrightarrow
\text{``}\varprojlim_n\text{''}\,\RR^{p+q}g_*(\mathcal J^n\overline K).
\]
Proposition~\ref{appen:prop5}(i) shows that this spectral sequence degenerates.\footnote{Editorial explanation: the terms of degree $p>0$ are zero pro-objects.} It then follows from
Proposition~\ref{appen:prop3} that the map obtained from
\eqref{appen:eq:independence} by passing to the formal inverse limit
is an isomorphism, giving the desired formula.

\subsection*{It remains to verify}

\begin{enumerate}
\item[(a)] a compatibility for every diagram. Denote the third compactification by $X'''$:
\[
\begin{tikzcd}[row sep=large,column sep=large]
X \ar[r,hook,"j'"] \ar[dr,hook,"j''"] \ar[ddr,hook,"j'''"'] \ar[ddd,"f"']
 & X' \ar[d,"h"] \\
 & X'' \ar[d,"g"] \\
 & X''' \ar[dl,"\overline f"] \\
S
\end{tikzcd}
\]
The successive identifications must be compatible:
\[
\begin{aligned}
\RR\overline f_*\RR j'''_!
&\cong\RR\overline f_*\RR g_*\RR j''_!\\
&\cong\RR\overline f_*\RR g_*\RR h_*\RR j'_!\\
&\cong\RR\overline f_*\RR(gh)_*\RR j'_!\\
&\cong\RR\overline f_*\RR j'''_!.
\end{aligned}
\]
Since two compactifications can always be dominated by a
third (for example $X'\times_S X''$),\footnote{Editorial clarification: take the scheme-theoretic closure of the diagonal image of $X$ in $X'\times_S X''$. The projections are proper and this closure contains $X$ as an open subset; the entire product does not in general contain $X$ as an open subset via this diagonal. See also \href{https://stacks.math.columbia.edu/tag/0ATU}{Stacks Project, Lemma 0ATU}. This clarification does not replace the compatibility verifications left outstanding.} this suffices to prove
independence of the choices.
\item[(b)] for every compactifiable pair, an identification
\[
\RR(fg)_!\cong\RR f_!\RR g_!.
\]
\item[(c)] for every compactifiable triple, a compatibility:
\[
\begin{tikzcd}[row sep=large,column sep=large]
\RR(fgh)_! \ar[r,"\sim"] \ar[d,"\sim"']
 & \RR(fg)_!\RR h_! \ar[d,"\sim"] \\
\RR f_!\RR(gh)_! \ar[r,"\sim"']
 & \RR f_!\RR g_!\RR h_!.
\end{tikzcd}
\]
\end{enumerate}

We would then have proved

\begin{booktheorem}{1}\label{appen:thm1}
For compactifiable $f:X\longrightarrow Y$, $K\in\operatorname{pro}\DD_{\operatorname{coh}}^b(X)$, and $L\in\operatorname{pro}\DD_{\operatorname{coh}}^b(Y)$, put $\Hom_f(K,L)=\Hom_Y(\RR f_!K,L)$. Subject to questions of compactifiability, this makes the categories $\operatorname{pro}\DD_{\operatorname{coh}}^b(X)$ into a category cofibred over noetherian preschemes (with compactifiable arrows, \dots).
\end{booktheorem}

\section*{No.\ 4. Definition of $\RR f^!$.}\label{appen:sec4}

It follows from general machinery of Verdier \hyperref[appen:bib1]{[1]} that the functor $\RR f_*:\DD(X)\longrightarrow\DD(S)$ has a \emph{right} adjoint $\DD^+(S)\longrightarrow\DD^+(X)$; it deserves a name, say $\RR f^!$, only when $f$ is proper, and its definition does not lend itself directly to calculation.

First one must make explicit a procedure for computing $\RR f_*$ at the level of complexes, of the form $\RR f_*K=f_*C^*(K)$, where $C^*$ is a finite acyclic resolution depending functorially and exactly on the object to which it is applied, and compatibly with filtered colimits. Then, for every injective $I$ on $S$, $\Hom(f_*C^p\mathcal F,I)$ is exact in the quasi-coherent sheaf $\mathcal F$ on $X$ and takes filtered $\varinjlim$ to filtered $\varprojlim$; it is therefore represented by a quasi-coherent injective $f_p^!I$. The $f_p^!I$ form a complex, defining $\RR f^!:\DD^+(S)\longrightarrow\DD^+(X)$, right adjoint to $\RR f_*:\DD(X)\longrightarrow\DD(S)$.

Here is how to define such a resolution:
\begin{enumerate}
\item[1)] If $S$ is separated: take a finite cover of $X$ by open subsets affine over $S$ (for example, affine open subsets), and set $\RR f_*K=f_*\check C(\mathfrak U,K)$ (the alternating \v{C}ech complex).

\item[2)] Otherwise: the canonical flasque resolutions have two defects here. (a) $f_*C^*$ is no longer quasi-coherent; this does not matter as long as, as here, quasi-coherent injectives are injective as sheaves. (b) $C^*$ does not commute with filtered colimits. This defect is easily corrected: express a quasi-coherent sheaf as $\mathcal F=\varinjlim\mathcal F_i$ (with the $\mathcal F_i$ finitely presented), and set $C'{}^*\mathcal F=\varinjlim C^*\mathcal F_i$. This is independent of the choices.
\end{enumerate}

For an open immersion, take $\RR f^!K=f^*K$. For a composite morphism $f=gh$ ($g$ proper, $h$ an open immersion), take $\RR f^!=\RR g^!\RR h^!$\footnote{Correction to the original: for $X\xrightarrow{h}\overline X\xrightarrow{g}S$, the contravariant functors compose in the order $\RR h^!\RR g^!$.}. Combining \bookref{appen:prop4}{Proposition 4} with the preceding discussion, we obtain

\begin{booktheorem}{2}\label{appen:thm2}
(Duality) Let $f:X\longrightarrow S$ be compactifiable, $f=gh$. The functors $\RR f_!:\operatorname{pro}\DD_{\operatorname{coh}}^b(X)\longrightarrow\operatorname{pro}\DD_{\operatorname{coh}}^b(S)$ and $\RR f^!:\DD^+(S)\longrightarrow\DD^+(X)$ are ``adjoint'' to one another.
\end{booktheorem}

One would have to verify independence of the compactification and a collection of compatibilities and identifications analogous to those for $\RR f_!$. Special cases of the formalism follow from the adjunction formula (uniqueness of an actual adjoint, \dots).

In the general case one can use the following proposition:

\begin{bookproposition}{6}\label{appen:prop6}
Let $f:X\longrightarrow S$ be compactifiable, and let $L\in\DD_{\operatorname{qc}}^+(S)$. The stalk at $x\in X$ of $\mathcal H^p\RR f^!L$ is given by
\[
(\mathcal{H}^{p}\RR f^{!}L)_{x} = \varinjlim_{x \in U} \Hom_{S}^{p}(\RR f_{!}(U,\mathcal{O}_{U}), L).
\]
\end{bookproposition}

We may assume that $X$ is proper and that $L$ is given as a complex of quasi-coherent injective sheaves. Let $U$ be an affine open subset of $X$, and let $\mathcal J$ be an ideal defining $X\setminus U$.
\begin{align*}
(\mathcal{H}^{p}\RR f^{!}L)(U)
 &= H^{p}\bigl((\RR f^{!}L)(U)\bigr) \\
 &= H^{p}\Hom_{\mathcal{O}_{U}}^{\bullet}
       (\mathcal{O}_{U},(\RR f^{!}L)|_{U}) \\
 &= H^{p}\varinjlim_{n}\Hom_{X}^{\bullet}
       (\mathcal{J}^{n},\RR f^{!}L) \\
 &= \varinjlim_{n}\Hom_{S}^{p}(\RR f_{*}\mathcal{J}^{n},L) \\
 &= \Hom_{S}^{p}(\RR f_{!}(U,\mathcal{O}_{U}),L).
\end{align*}
The proposition follows.

We shall assume for $\RR f^!$ a variance formalism analogous to that of \bookref{appen:thm1}{Theorem 1} for $\RR f_!$.

To make the definition of $\RR f^!$ in the proper case more explicit, one may observe that if a quasi-coherent sheaf $\mathcal F$ represents a functor $F$, then for every open subset $U$ and every ideal $\mathcal J$ defining $X\setminus U$, one has $\mathcal F(U)=\varinjlim\Hom(\mathcal J^n,\mathcal F)=\varinjlim F(\mathcal J^n)$.

\section*{No.\ 5. Calculation of $\RR f^!$.}\label{appen:sec5}

For a finite morphism, $\RR f^!$ is what we want: to see this, it suffices to use the identity as the resolving functor in the calculation of $\RR f_*$. For projective space, uniqueness of an adjoint leaves no choice (using the fact that we have the adjunction formula for $\RR f_!=\RR f_*:\DD(X)\to\DD(S)$); we shall return to the map. For an open immersion, $\RR f^!=\RR f^*$. Thus, for every open subset $U$ of $X$ that is quasi-projective over $S$, one can calculate the restriction of $\RR f^!$ to $U$. In particular, if one can verify locally that $\mathcal H^q\RR f^!K=0$ for $q\ne p$, then $\RR f^!K$ is known, since an object of the derived category having just one nonzero cohomology sheaf is determined by that sheaf. Recall an important case in which this condition holds, with $K=\mathcal O$.

\begin{samepage}
\begin{bookproposition}{7}\label{appen:prop7}
Let $f:X\to Y$ be compactifiable. Suppose that locally (on $X$ and $Y$) one can find a diagram
\[
\begin{tikzcd}[column sep=large,row sep=large]
X \ar[r,"f"] \ar[dr,"g"'] & Y \ar[d,"h"] \\
 & S
\end{tikzcd}
\]
in which $g$ and $h$ are local complete intersections of relative dimensions $d'$ and $d''$, and put $d=d'-d''$. Then $\RR f^!\mathcal O$ has only one cohomology sheaf, in degree $-d$, and that sheaf is invertible.
\end{bookproposition}
\end{samepage}

Indeed, $\RR g^!\mathcal O=\RR f^!\RR h^!\mathcal O$; we know that $\RR g^!\mathcal O$ and $\RR h^!\mathcal O$ are of the indicated form, in degrees $-d'$ and $-d''$, by local calculations in projective space. Since everything is local, we may assume that $\mathcal H^{-d''}\RR h^!\mathcal O$ is isomorphic to $\mathcal O$, and the assertion follows.

To pass from this to more general complexes, one would first have to prove formulas of the form $\RR f^!(K\derivedtensor L)=\RR f^!(K)\derivedtensor\LL f^*L$. Only the definition of the map presents a problem, since being an isomorphism is a local question. Restrictions on degrees are of course necessary. I have not checked any of what follows in detail.

\subsection*{First method for defining the map.}

Suppose that $K$ and $L$ belong to $\DD_{\operatorname{coh}}^b(S)$ (and are not merely quasi-coherent), and suppose in addition that
\begin{enumerate}
\item[(a)] $\RR f^!K$ is bounded in degree (it is automatically coherent, since this is a local question);
\item[(b)] $K\derivedtensor L$ is bounded in degree;
\item[(c)] $\RR f^!K\derivedtensor\LL f^*L$ is bounded in degree.
\end{enumerate}

To define
\[
\RR f^{!}K \derivedtensor \LL f^{*}L
 \longrightarrow \RR f^{!}(K \derivedtensor L),
\]
it suffices, by duality, to define
\[
\RR f_{!}(\RR f^{!}K \derivedtensor \LL f^{*}L)
 \longrightarrow K \derivedtensor L,
\]
and by K\"unneth
\[
\RR f_{!}(\RR f^{!}K \derivedtensor \LL f^{*}L)
 = \RR f_{!}\RR f^{!}K \derivedtensor L;
\]
the desired map is induced by the adjunction map
$\RR f_{!}\RR f^{!}K \longrightarrow K$.

\subsection*{Second method.}

Suppose that $f$ is compactifiable, flat, and a relative local complete intersection. Then $\RR f^!\mathcal O=\omega[d]$, with $\omega$ invertible. Let $\RR f_!$ denote a procedure for computing $\RR f_!$ at the level of complexes (that is, such a procedure for a compactification $\overline f:\overline X\to S$ of $f$), assumed to be of the type described in \bookref{appen:sec4}{No.\ 4}, and such that $C^p\mathcal F$ vanishes for $p>d$ (which can be achieved by truncation). There is then a map
\[
R^{d}\overline{f}_{*}\mathcal{F}[-d]
 \longrightarrow \RR\overline{f}_{*}\mathcal{F}
\]
for every $\mathcal F$ on $\overline X$\footnote{Editorial correction and qualification: for the calculation on the compactification, read $\overline f_*$ in the Hom complexes below. If $C=\RR\overline f_*\mathcal F$ has cohomology in degrees $\le d$, canonical truncation gives $C\to R^d\overline f_*\mathcal F[-d]$, in the direction opposite to the printed arrow; see \href{https://stacks.math.columbia.edu/tag/08J5}{Stacks Project, Remark 08J5}. Applying $\Hom_S^{\bullet}(-,I)$ reverses the direction once more: the comparison thus obtained is $\RR f_1^!I\to\RR f^!I$.
\par The required bound is $R^q\overline f_*\mathcal F=0$ for $q>d$, for the sheaves under consideration on $\overline X$. It does not follow merely from the relative dimension of the open morphism $f$. For example, let $S=\mathbb A_k^2$ and $f:X=S\setminus\{0\}\hookrightarrow S$, of relative dimension $d=0$. The compactification $\overline f:\operatorname{Bl}_0 S\to S$ has an exceptional divisor $i:E\cong\mathbb P_k^1\hookrightarrow\overline X$; for $\mathcal F=i_*\sO_E(-2)$, one has $R^1\overline f_*\mathcal F\cong k(0)\ne0$. Thus this model cannot be truncated uniformly in degree $d=0$. This example concerns the procedure on an arbitrary compactification; it does not contradict the final formula for $f^!$.} (the last cohomology sheaf maps into its complex). This gives morphisms of complexes
\[
\Hom_{S}^{\bullet}(\RR f_{*}\mathcal{F}, I)
 \longrightarrow
\Hom_{S}^{\bullet}(R^{d}f_{*}\mathcal{F}[-d], I)
\]
for every complex $I$ on $S$. Let $I$ be injective. Then $R^d f_*\mathcal F[-d]$ is right exact in $\mathcal F$, so the functor in $\mathcal F$ on the right of the arrow is representable: define $I_1$ by
\[
\Hom_{S}(R^{d}f_{*}\mathcal{F}, I)
 = \Hom_{\overline X}(\mathcal{F}, I_1),
 \qquad \RR f_{1}^{!}I = I_1[d].
\]
The preceding construction defines
\[
\RR f^{!}I \longrightarrow \RR f_{1}^{!}I
\]
and for a complex of injectives, $\RR f_1^!$ is defined term by term. It remains to evaluate $\RR f_1^!I$ on affine open subsets $U$ of $X$; this is given by $\varinjlim_n\Hom(R^d f_!\mathcal J^n,I)$. Since this problem is local on the source, it is easy to show by projective methods that the dual of
\[
\text{``}\varprojlim_{n}\text{''}\,R^{d}f_{!}\mathcal{J}^{n}
 = R^{d}f_{!}(U,\mathcal{O}),
\]
is the Ind-object\footnote{Editorial clarification: for $P=\text{``}\varprojlim_n\text{''}P_n$ and an injective $I$, one uses $\Hom_S(P,I)=\varinjlim_n\Hom_S(P_n,I)$. Passing to Hom therefore turns the inverse system into a direct system. Identifying this colimit with $\Gamma(U,\omega\otimes f^*I)$, compatibly with restriction in $U$, is the local calculation sketched by the author; it is not justified by a simple termwise interchange of Hom and tensor. The shift $[d]$ places the sheaf in degree $-d$. Since $f$ is flat and $\omega$ is invertible, $\omega\otimes f^*(-)$ is exact; once this identification has been obtained, the termwise calculation respects quasi-isomorphisms. The compatibilities announced but not verified remain unverified.} $f_*(U,\omega)$, represented as the colimit of its submodules of sections having a pole of order at most $n$ outside $U$ ($n\to\infty$).

This gives $\RR f_1^!I=\omega\otimes f^*I[d]$, and the general formula $\RR f_1^!K=f^*K\otimes\omega[d]$.

\section*{Bibliography (for the Appendix)}
\begin{enumerate}
\item[{[1]}]\phantomsection\label{appen:bib1} J.L.\ Verdier, S\'eminaire Bourbaki, November 1965, expos\'e 300.
\item[{[2]}]\phantomsection\label{appen:bib2} M.\ Nagata, Imbedding of an abstract variety in a complete variety. Journal of Math.\ of Kyoto University, vol.\ 2, no.\ 1, 1962.
\end{enumerate}
\par\endgroup
